(1) 输入慢变复包络为 \(A(0,t)=A_0\exp[-(1+iC)t^2/(2\tau_0^2)]\)，其中 \(C\) 为无量纲初始线性啁啾，\(\tau_0>0\)。传播长度 \(L\) 的纯二阶色散介质在频域乘 \(e^{i\beta_2L\Omega^2/2}\)（时间傅里叶因子 \(e^{-i\Omega t}\)）。求传播后的Gaussian复宽度。

(2) 求输出强度的时间宽度 \(\tau_L\) 与峰值功率比，证明无损色散不改变脉冲总能量。

(3) 优化可以自由选取的色散量 \(\beta_2L\) 使脉冲最短，求最佳色散、最短宽度及峰值增强。

(4) 讨论无初始啁啾和色散符号固定时的极限，并解释为什么频域二次相位不会改变功率谱。
